
\section{Dual Prediction and Feasibility Restoration}
\label{sec:method}

We now construct a prediction rule $\pi_{\boldsymbol\theta}$
that combines a learned mapping for predicting subproblem
dual solutions (Section~\ref{subsec:direct}) with a
feasibility-restoration procedure that avoids solving
an auxiliary optimization problem
(Section~\ref{subsec:repair}).

\subsection{State-Aware Dual Learning}
\label{subsec:direct}

Section~\ref{subsec:labd} identifies three sources of
information for predicting the dual solution of
$m\in\mathcal M_p$: its fixed representation
$\mathbf z_m$, its current dual-objective state
$\mathbf g_m(\mathbf x)$, and the current exact duals
$\boldsymbol\lambda_e(\mathbf x)$. The learned mapping
uses the first two to capture persistent subproblem
characteristics and the current state, while the exact
duals provide a feasible reference point for the
restoration step. Define
\begin{equation}\label{eq:state-transformation}
\mathbf s_m(\mathbf x):=
\phi_g\bigl(\mathbf g_m(\mathbf x)\bigr)\in\R^q,
\end{equation}
where $\phi_g$ is the identity map or a fixed
transformation of the dual-objective state.
Consider a finite-dimensional parametric mapping
$f_{\boldsymbol\theta}:\mathcal Z\times\R^q\to\R^r$.
The raw prediction for subproblem $m$ at $\mathbf x$ is
\begin{equation}\label{eq:raw-prediction}
\widetilde{\boldsymbol\lambda}_m(\mathbf x)
:=f_{\boldsymbol\theta}
\bigl(\mathbf z_m,\mathbf s_m(\mathbf x)\bigr).
\end{equation}
The warm-up panel~\eqref{eq:panel} provides the exact
dual responses for training. With
$\mathbf s_m^k:=\mathbf s_m(\mathbf x^k)$,
a general empirical learning formulation is
\begin{equation}\label{model:training}
\min_{\boldsymbol\theta\in\Theta}\quad
\frac{1}{|\mathcal K_0||\mathcal M|}
\sum_{k\in\mathcal K_0}\sum_{m\in\mathcal M}
\ell\Bigl(
\boldsymbol\lambda_m^k,\,
f_{\boldsymbol\theta}(\mathbf z_m,\mathbf s_m^k);
\mathbf g_m^k
\Bigr)+\mathcal R(\boldsymbol\theta),
\end{equation}
where $\ell$ is a nonnegative prediction loss and
$\mathcal R$ is a regularization function.
After training, the selected set $\mathcal M_e$,
representations $\mathbf z_m$, transformation $\phi_g$,
and fitted parameters $\boldsymbol\theta^*$ remain fixed;
subsequent predictions use the current state.

\subsection{Dual-Feasibility Restoration}
\label{subsec:repair}

The predicted dual vector
$\widetilde{\boldsymbol\lambda}_m(\mathbf x)\in\R^r$
may lie outside the dual-feasible polyhedron
$\mathcal D$. A natural approach is to project it
back onto $\mathcal D$, for example through the
$\ell_1$-projection
\begin{equation}\label{eq:projection}
\min_{\boldsymbol\lambda\in\mathcal D}
\bigl\|\boldsymbol\lambda-
\widetilde{\boldsymbol\lambda}_m(\mathbf x)\bigr\|_1.
\end{equation}
This requires solving an auxiliary linear program
for each predicted subproblem at every cut-generating
solution. Moreover, proximity to the prediction
does not necessarily imply a stronger Benders cut.

We instead exploit the polyhedral structure of
$\mathcal D$ to perform a directional projection
along the segment joining the prediction to a known
feasible point $\bar{\boldsymbol\lambda}\in\mathcal D$.
The maximum feasible step is determined explicitly
by the defining inequalities of $\mathcal D$.
Linearity of the dual objective then allows us to
select the step that maximizes the cut value at
the generating solution.

\begin{proposition}[Maximum feasible step]
\label{prop:repair}
Let $\bar{\boldsymbol\lambda}\in\mathcal D$,
$\widetilde{\boldsymbol\lambda}\in\R^r$, and
$\mathbf d:=\widetilde{\boldsymbol\lambda}
-\bar{\boldsymbol\lambda}$. Then
\[
\{\gamma\in[0,1]:
\bar{\boldsymbol\lambda}+\gamma\mathbf d\in\mathcal D\}
=[0,\gamma^{\max}],
\]
where
\begin{equation}\label{eq:gamma-max}
\gamma^{\max}
=\min\Bigl\{
1,\
\min_{j:\,d_j<0}\frac{\bar\lambda_j}{-d_j},\
\min_{\ell:\,(\mathbf W^\top\mathbf d)_\ell<0}
\frac{(\mathbf W^\top\bar{\boldsymbol\lambda})_\ell+q_\ell}
{-(\mathbf W^\top\mathbf d)_\ell}
\Bigr\},
\end{equation}
with minima over empty index sets omitted.
Moreover, for any $\mathbf g\in\R^r$, the dual
objective $(\bar{\boldsymbol\lambda}+
\gamma\mathbf d)^\top\mathbf g$ is maximized
over $[0,\gamma^{\max}]$ at
\begin{equation}\label{eq:guard}
\gamma^\star=
\begin{cases}
\gamma^{\max}, & \mathbf d^\top\mathbf g>0,\\
0, & \text{otherwise}.
\end{cases}
\end{equation}
\end{proposition}

Proposition~\ref{prop:repair} replaces the auxiliary
projection problem with an explicit feasibility
calculation along a line segment. The rule
\eqref{eq:guard} further ensures that the resulting
cut is no weaker than the reference cut at the
generating solution.

To obtain a feasible reference point, we exploit
the convexity of the common dual-feasible region
$\mathcal D$. The average of the exact dual solutions
at the current first-stage solution,
\begin{equation}\label{eq:mean-anchor}
\bar{\boldsymbol\lambda}(\mathbf x)
:=\frac{1}{p}\sum_{t\in\mathcal M_e}
\boldsymbol\lambda_t(\mathbf x),
\end{equation}
belongs to $\mathcal D$ and requires no additional
subproblem evaluations. Applying
Proposition~\ref{prop:repair} with this reference
point and
$\mathbf d_m(\mathbf x):=
\widetilde{\boldsymbol\lambda}_m(\mathbf x)
-\bar{\boldsymbol\lambda}(\mathbf x)$
yields the restored dual prediction
\begin{equation}\label{eq:certified}
\widehat{\boldsymbol\lambda}_m(\mathbf x)
:=
\bigl(1-\gamma_m(\mathbf x)\bigr)
\bar{\boldsymbol\lambda}(\mathbf x)
+\gamma_m(\mathbf x)
\widetilde{\boldsymbol\lambda}_m(\mathbf x),
\end{equation}
where $\gamma_m(\mathbf x)$ is defined by
\eqref{eq:guard} with
$\mathbf g=\mathbf g_m(\mathbf x)$.
This defines the prediction rule
$\pi_{\boldsymbol\theta^*}
(\mathbf z_m;\mathbf g_m(\mathbf x),
\boldsymbol\lambda_e(\mathbf x))
:=\widehat{\boldsymbol\lambda}_m(\mathbf x)$.

\begin{proposition}[Feasibility and cut improvement]
\label{prop:certificate}
For every $\mathbf x\in\mathcal X$ and
$m\in\mathcal M_p$,
$\widehat{\boldsymbol\lambda}_m(\mathbf x)
\in\mathcal D$, and
\begin{equation}\label{eq:improvement}
\begin{aligned}
\widehat{\boldsymbol\lambda}_m(\mathbf x)^\top
\mathbf g_m(\mathbf x)
&=\bar{\boldsymbol\lambda}(\mathbf x)^\top
\mathbf g_m(\mathbf x)\\
&\quad+
\gamma_m^{\max}(\mathbf x)
\bigl[\mathbf d_m(\mathbf x)^\top
\mathbf g_m(\mathbf x)\bigr]_+\\
&\ge
\bar{\boldsymbol\lambda}(\mathbf x)^\top
\mathbf g_m(\mathbf x),
\end{aligned}
\end{equation}
where $[u]_+:=\max\{u,0\}$.
\end{proposition}
